% WWW 2027 research track, Security, Privacy and Trust. Check CfP for page limit & template options.
\documentclass[sigconf,anonymous,review]{acmart}
\usepackage{booktabs}
\newcommand{\todo}[1]{\textcolor{red}{[TODO: #1]}}
\newcommand{\verify}[1]{\textcolor{orange}{#1}} % numbers not yet re-verified (see AUDIT.md)

\begin{document}
\title{What Does an XSS Detector Actually Detect?\\An Execution-Grounded Measurement of Learned Cross-Site Scripting Detection}

\begin{abstract}
Machine-learning detectors for cross-site scripting (XSS) routinely report accuracies above 99\%.
These figures are computed against a label meaning ``a dataset called this string an attack''.
Operators care about something else: whether a payload \emph{executes}.
We build a local testbed (browser, inline proxy hosting the detector, target applications,
four production template engines) and label payloads by whether an inert marker fires in a
real browser at each sink.
Across ten detector architectures, scores that separate attack from benign almost perfectly
are \emph{anti}-correlated with execution (\verify{0.34--0.51} AUC, grouped by payload family).
A length-only baseline explains much of this: long, exotic cheat-sheet strings look most
attack-like and run least.
Public XSS corpora are largely enumerations of tag--event pairs: only \verify{2--3\%} of payloads
execute on render, and the rest need a user action, a precondition or a legacy browser,
for reasons we categorise.
Among sinks where execution is possible, the output context adds \verify{+0.08} AUC over the
payload, a gap no request-side detector can close.
Finally, synthetic event dispatch, the usual shortcut for verifying interaction-gated
payloads, reaches elements no user can, and over-reports exploitability.
We release the testbed, oracle and labels.
\end{abstract}

\maketitle

\section{Introduction}
% Para 1: the 99% numbers, and what label they are computed against.
% Para 2: why execution is the label that matters (AUDIT M3):
%   alert triage; evasion/robustness papers that never check the evading payload still runs;
%   benchmark inflation from dead payloads.
% Para 3: our approach: execution-grounded labels, two labels never mixed.
% Para 4: no new model; a measurement. Pitfalls framing, cf. \cite{arp2022dos}.
\paragraph{Contributions}
\begin{itemize}
  \item An execution oracle and testbed that keep \emph{attack-likeness} and \emph{executability} separate, with the oracle's own error quantified in both directions (\S\ref{sec:oracle}).
  \item RQ1: label validity across ten detector architectures and two published detectors (\S\ref{sec:rq1}).
  \item RQ2: the role of output context, incl.\ context-mismatched escaping in real template engines (\S\ref{sec:rq2}).
  \item RQ3: inspection-point coverage for reflected / stored / DOM XSS (\S\ref{sec:rq3}).
  \item RQ4: decay of public payload benchmarks, with a taxonomy of why payloads are dead (\S\ref{sec:rq4}).
\end{itemize}

\section{Background and Related Work}
% ML XSS detection (DeepXSS \cite{fang2018deepxss}, Miczek et al. \cite{miczek2025llm}); WAFs and evasion (\cite{demetrio2020wafamole});
% execution-verified XSS discovery (\cite{lekies2013flows,melicher2018domsday,steffens2019locals,duchene2014kameleonfuzz});
% ML-in-security pitfalls (\cite{arp2022dos}). Position: they verify execution to *find* bugs; we use it to *audit detectors*.

\section{Methodology}
\subsection{Two labels}
\subsection{Testbed}            % three processes; why the boundary is the experiment
\subsection{Payload corpus}     % sources, dedup by skeleton, group-level split, licences
\subsection{Execution oracle}\label{sec:oracle}
% marker + instrumented hooks (AUDIT B1); render-settle protocol (B2);
% synthetic vs real input, BOTH directions (B3); cross-browser kappa (M10).
\subsection{Statistics}         % cluster bootstrap over group_id (B4)

\section{RQ1: Do detectors track execution?}\label{sec:rq1}
% Table 3 (our 10 archs) + CRS PL1–4 real engine (B6) + Miczek-style RF; positive control trained on execution (M4).
% Split sensitivity: report ranges.

\section{RQ2: Where the value lands}\label{sec:rq2}
% Per-sink rates; escaped text = positive control (M2); context-mismatched escaping = the result.
% Payload vs sink vs both; escaped sinks excluded (M5); leave-one-sink-out (B5).

\section{RQ3: Which side of the wire}\label{sec:rq3}
% request / response / response+XHR per class (M6 wording).

\section{RQ4: Benchmark decay}\label{sec:rq4}
% Unmodified public rows; taxonomy (M9); do LLM-obfuscated payloads still execute? (M3b)

\section{Discussion}
% Implications: for detector authors (report execution-grounded metrics), for benchmark
% curators, for WAF operators (context-aware / response-side placement).
% Negative results: augmentation, calibration, cross-app transfer.

\section{Limitations}
\section{Ethics}
% All local, 127.0.0.1, inert marker, no third-party systems tested.

\section{Conclusion}

\bibliographystyle{ACM-Reference-Format}
\bibliography{refs}
\end{document}
